\Needspace{18\baselineskip}
\subsection{A polar metric in a background with unbounded mass}
\label{subsec:unbounded-auxiliary-polar-metric}

Taking the Hopf threefold of \cref{thm:hopf-unbounded} as the auxiliary factor makes the ambient masses unbounded, while the complete metric on the product models is unchanged. Letting the coefficient of the Hopf factor vary over the curve product also makes the background nonclosed.

\Needspace{18\baselineskip}
\begin{example}\label{ex:unbounded-auxiliary-polar-metric}\index{Metric completeness!unbounded ambient mass}\index{Bounded mass property!complete zero-volume model locus}
There is a compact complex sixfold with a smooth positive background $\Theta$ that is not closed and has infinite upper mass $v_+(\Theta)=+\infty$ in the sense of \cref{def:bmp}. Nevertheless, the locus of separated product models of \cref{prop:general-auxiliary-product-metric}, along a divisor of weight one, has a complete geodesic polar metric of diameter $4$. The ideal subspace is closed and is the metric closure of the models with analytic singularities, and ideal projection is the unique nearest-point retraction onto it and is $1$-Lipschitz. The upper-volume comparison vanishes on this locus. The locus contains distinct ideal models with analytic singularities at polar distance $1$.
\end{example}
\begin{proof}
Take $Y=C_1\times C_2$ with $C_1=C_2=\mathbb P^1$ and normalized Fubini--Study forms $\omega_i$ of area one, and write $\pi_i\colon Y\to C_i$ for the projections. Let $(T,\omega_T,V_j)$ be the Hopf threefold, its Gauduchon form and the smooth potentials of \cref{thm:hopf-unbounded}. Take $Z=\mathbb P^1$ with its normalized Fubini--Study form $\kappa$, a point $D$ and its normalized potential $\psi$ satisfying $\kappa+\ddc\psi=[D]$. On $M=Y\times T\times Z$, write $p,s,r$ for the projections. Choose a positive nonconstant smooth function $h$ on $Y$ and put
\[
 \Theta=p^*(\pi_1^*\omega_1+\pi_2^*\omega_2)
          +(h\circ p)s^*\omega_T+r^*\kappa,
 \qquad \Psi=\psi\circ r.
\]
This form is strictly positive, and it is not closed by \cref{cor:general-auxiliary-background-boundaries}.

Put $c=\min_Yh>0$. The smooth functions $W_j=cV_j\circ s$ are strictly $\Theta$-plurisubharmonic, since their current in the $T$-directions is
\[
 h\omega_T+c\ddc V_j
    =c(\omega_T+\ddc V_j)+(h-c)\omega_T>0.
\]
\Needspace{6\baselineskip}
Only the part of degree three in the $T$-directions contributes to the top-degree product. The three curve forms have integral one, so positivity of smooth mixed products and Fubini's theorem give
\[
 \begin{aligned}
 \int_M(\Theta+\ddc W_j)^6
 &=120\int_Y\pi_1^*\omega_1\wedge\pi_2^*\omega_2
                  \int_T(h(y)\omega_T+c\ddc V_j)^3\\
 &\geq120c^3\int_T(\omega_T+\ddc V_j)^3
       \longrightarrow+\infty.
 \end{aligned}
\]
The factor $120=6!/3!$ counts the ways of choosing one factor for each curve direction and three for $T$. Each $W_j$ is bounded, so \cref{def:bmp} gives $v_+(\Theta)=+\infty$, and the bounded mass property fails.

Apply \cref{prop:general-auxiliary-product-metric} to this product, with both curve areas equal to one. On the locus of potentials $\Psi+q_1\circ\pi_1\circ p+q_2\circ\pi_2\circ p$, where each $q_i$ ranges over the ordinary $\omega_i$-models with supremum zero, it gives a complete metric $\rho$ of diameter $4$, the analytic closure of the ideal subspace, the nearest-point retraction and the vanishing of the upper-volume comparison.

For $x\in C_1$, let $g_x$ be the normalized Green potential of the point $x$, with $\omega_1+\ddc g_x=[x]$. By \cref{lem:general-auxiliary-product-projections}, the potentials
\[
 \Phi_0=\Psi,\qquad
 \Phi_x=\Psi+g_x\circ\pi_1\circ p
\]
are ordinary and ideal models, and their divisorial poles with coefficient one give analytic singularities. Their tuples of polar curvatures are $(0,0)$ and $(\delta_x,0)$, so $\rho(\Phi_0,\Phi_x)=\|\delta_x\|_{\mathrm{TV}}=1$. Nevertheless, by the same lemma, every representative of either ordinary type has zero extended mass; in particular, all their upper and lower singularity volumes are zero.
\end{proof}
